\chapter{2023年天津卷}

\section{单选题}

\begin{problem}
已知集合 \(U=\{1,2,3,4,5\}\),\(A=\{1,3\}\),\(B=\{1,2,4\}\),则 \(\complement_U B\cup A=\)(\quad)

\choices
    {\(\{1,3,5\}\)}
    {\(\{1,3\}\)}
    {\(\{1,2,4\}\)}
    {\(\{1,2,4,5\}\)}
\end{problem}

\begin{problem}
``\(a^2=b^2\)''是``\(a^2+b^2=2ab\)''的(\quad)

\choices
    {充分不必要条件}
    {必要不充分条件}
    {充分必要条件}
    {既不充分又不必要条件}
\end{problem}

\begin{problem}
若 \(a=1.01^{0.5}\),\(b=1.01^{0.6}\),\(c=0.6^{0.5}\),则 \(a\),\(b\),\(c\) 的大小关系为(\quad)

\choices
    {\(c>a>b\)}
    {\(c>b>a\)}
    {\(a>b>c\)}
    {\(b>a>c\)}
\end{problem}

\begin{problem}
函数 \(f(x)\) 的图象如下图所示,则 \(f(x)\) 的解析式可能为(\quad)

\bitmapfigure[max width=0.62\linewidth]{img/2023/tianjin/q4_fig1.png}

\choices
    {\(\dfrac{5(\e^x-\e^{-x})}{x^2+2}\)}
    {\(\dfrac{5\sin x}{x^2+1}\)}
    {\(\dfrac{5(\e^x+\e^{-x})}{x^2+2}\)}
    {\(\dfrac{5\cos x}{x^2+1}\)}
\end{problem}

\begin{problem}
已知函数 \(f(x)\) 的一条对称轴为直线 \(x=2\),一个周期为 4,则 \(f(x)\) 的解析式可能为(\quad)

\choices
    {\(\sin \left( \dfrac{\pi}{2}x \right)\)}
    {\(\cos \left( \dfrac{\pi}{2}x \right)\)}
    {\(\sin \left( \dfrac{\pi}{4}x \right)\)}
    {\(\cos \left( \dfrac{\pi}{4}x \right)\)}
\end{problem}

\begin{problem}
已知 \(\{a_n\}\) 为等比数列,\(S_n\) 为数列 \(\{a_n\}\) 的前 \(n\) 项和,\(a_{n+1}=2S_n+2\),则 \(a_4=\)(\quad)

\choices
    {3}
    {18}
    {54}
    {152}
\end{problem}

\begin{problem}
调查某种群花萼长度和花瓣长度,所得数据如图所示,其中相关系数 \(r=0.8245\),下列说法正确的是(\quad)

\examfiguregroup[float=inline]{img/2023/tianjin/q7_fig1.png}{%
    \begin{tabular}{cc}
    \bitmapinclude[max width=6.0cm]{img/2023/tianjin/q7_fig1.png} &
    \bitmapinclude[max width=5.0cm]{img/2023/tianjin/q7_fig2.png}
    \end{tabular}%
}

\choices
    {花瓣长度和花萼长度没有相关性}
    {花瓣长度和花萼长度呈现负相关}
    {花瓣长度和花萼长度呈现正相关}
    {若从样本中抽取一部分,则这部分的相关系数一定是 0.8245}
\end{problem}

\begin{problem}
在三棱锥 \(P-ABC\) 中,线段 \(PC\) 上的点 \(M\) 满足 \(PM=\dfrac{1}{3}PC\),线段 \(PB\) 上的点 \(N\) 满足 \(PN=\dfrac{2}{3}PB\),则三棱锥 \(P-AMN\) 和三棱锥 \(P-ABC\) 的体积之比为(\quad)

\choices
    {\(\dfrac{1}{9}\)}
    {\(\dfrac{2}{9}\)}
    {\(\dfrac{1}{3}\)}
    {\(\dfrac{4}{9}\)}
\end{problem}

\begin{problem}
双曲线 \(\frac{x^2}{a^2}-\frac{y^2}{b^2}=1\)(\(a>0\),\(b>0\)) 的左,右焦点分别为 \(F_1\),\(F_2\). 过 \(F_2\) 作其中一条渐近线的垂线,垂足为 \(P\). 已知 \(PF_2=2\),直线 \(PF_1\) 的斜率为 \(\dfrac{\sqrt{2}}{4}\),则双曲线的方程为(\quad)

\bitmapfigure[max width=0.42\linewidth]{img/2023/tianjin/q9_fig1.png}

\choices
    {\(\dfrac{x^2}{8}-\dfrac{y^2}{4}=1\)}
    {\(\dfrac{x^2}{4}-\dfrac{y^2}{8}=1\)}
    {\(\dfrac{x^2}{4}-\dfrac{y^2}{2}=1\)}
    {\(\dfrac{x^2}{2}-\dfrac{y^2}{4}=1\)}
\end{problem}

\section{填空题}

\begin{problem}
已知 \(\i\) 是虚数单位,化简 \(\frac{5+14\i}{2+3\i}\) 的结果为\fillinblank{}.
\end{problem}

\begin{problem}
在 \(\left( 2x^3-\frac{1}{x} \right)^6\) 的展开式中,\(x^2\) 项的系数为\fillinblank{}.
\end{problem}

\begin{problem}
过原点的一条直线与圆 \(C:(x+2)^2+y^2=3\) 相切,交曲线 \(y^2=2px\)(\(p>0\)) 于点 \(P\),若 \(|OP|=8\),则 \(p\) 的值为\fillinblank{}.
\end{problem}

\begin{problem}
甲,乙,丙三个盒子中装有一定数量的黑球和白球,其总数之比为 \(5:4:6\). 这三个盒子中黑球占总数的比例分别为 \(40\%\),\(25\%\),\(50\%\). 现从三个盒子中各取一个球,取到的三个球都是黑球的概率为\fillinblank{};将三个盒子混合后任取一个球,是白球的概率为\fillinblank{}.
\end{problem}

\begin{problem}
在 \(\triangle ABC\) 中,\(\angle A=60^\circ\),\(BC=1\),点 \(D\) 为 \(AB\) 的中点,点 \(E\) 为 \(CD\) 的中点. 若设 \(\overrightarrow{AB}=\bs{a}\),\(\overrightarrow{AC}=\bs{b}\),则 \(\overrightarrow{AE}\) 可用 \(\bs{a}\),\(\bs{b}\) 表示为\fillinblank{};若 \(\overrightarrow{BF}=\dfrac{1}{3}\overrightarrow{BC}\),则 \(\overrightarrow{AE}\cdot \overrightarrow{AF}\) 的最大值为\fillinblank{}.
\end{problem}

\begin{problem}
若函数 \(f(x)=ax^2-2x-\left| x^2-ax+1 \right|\) 有且仅有两个零点,则 \(a\) 的取值范围为\fillinblank{}.
\end{problem}

\section{解答题}

\begin{problem}
在 \(\triangle ABC\) 中,角 \(A\),\(B\),\(C\) 所对的边分别是 \(a\),\(b\),\(c\). 已知 \(a=\sqrt{39}\),\(b=2\),\(\angle A=120^\circ\).

\begin{enumerate}
\item 求 \(\sin B\) 的值;
\item 求 \(c\) 的值;
\item 求 \(\sin(B-C)\).
\end{enumerate}
\end{problem}

\begin{problem}
三棱台 \(ABC-A_1B_1C_1\) 中,若 \(A_1A\perp \text{面 }ABC\),\(AB\perp AC\),\(AB=AC=A_1A=2\),\(A_1C_1=1\),\(M\),\(N\) 分别是 \(BC\),\(BA\) 中点.

\begin{enumerate}
\item 求证:\(A_1N\parallel \text{平面 }C_1MA\);
\item 求平面 \(C_1MA\) 与平面 \(ACC_1A_1\) 所成夹角的余弦值;
\item 求点 \(C\) 到平面 \(C_1MA\) 的距离.
\end{enumerate}

\bitmapfigure[max width=0.52\linewidth]{img/2023/tianjin/q17_fig1.png}
\end{problem}

\begin{problem}
设椭圆 \(\frac{x^2}{a^2}+\frac{y^2}{b^2}=1\)(\(a>b>0\)) 的左,右顶点分别为 \(A_1\),\(A_2\),右焦点为 \(F\),已知 \(|A_1F|=3\),\(|A_2F|=1\).

\begin{enumerate}
\item 求椭圆方程及其离心率;
\item 已知点 \(P\) 是椭圆上一动点(不与端点重合),直线 \(A_2P\) 交 \(y\) 轴于点 \(Q\),若三角形 \(A_1PQ\) 的面积是三角形 \(A_2FP\) 面积的二倍,求直线 \(A_2P\) 的方程.
\end{enumerate}
\end{problem}

\begin{problem}
已知 \(\{a_n\}\) 是等差数列,\(a_2+a_5=16\),\(a_5-a_3=4\).

\begin{enumerate}
\item 求 \(\{a_n\}\) 的通项公式和 \(\displaystyle\sum_{i=2^{n-1}}^{2^n-1}a_i\);
\item 已知 \(\{b_n\}\) 为等比数列,对于任意 \(k\in \N^*\),若 \(2^{k-1}\le n\le 2^k-1\),则 \(b_k<a_n<b_{k+1}\).

    \begin{enumerate}
    \item 当 \(k\ge 2\) 时,求证:\(2^k-1<b_k<2^k+1\);
    \item 求 \(\{b_n\}\) 的通项公式及其前 \(n\) 项和.
    \end{enumerate}
\end{enumerate}
\end{problem}

\begin{problem}
已知函数 \(f(x)=\left( \frac{1}{x}+\frac{1}{2} \right)\ln(x+1)\).

\begin{enumerate}
\item 求曲线 \(y=f(x)\) 在 \(x=2\) 处切线的斜率;
\item 当 \(x>0\) 时,证明:\(f(x)>1\);
\item 证明:
\(\frac{5}{6}<\ln(n!)-\left( n+\frac{1}{2} \right)\ln n+n\le 1\)
\end{enumerate}
\end{problem}

